\documentclass[11pt]{article}
\usepackage[a4paper,margin=18mm]{geometry}
\usepackage{fontspec}
\setmainfont{Latin Modern Roman}
\usepackage{amsmath,amssymb,amsthm,amscd,graphicx,color}
\usepackage{xurl}
\usepackage[unicode,hidelinks]{hyperref}
\allowdisplaybreaks
\setlength\emergencystretch{3em}

\makeatletter
\newtheorem{theorem}{Theorem}
\newtheorem*{theorem*}{Theorem}
\newtheorem{lemma}{Lemma}
\newtheorem*{lemma*}{Lemma}
\newtheorem{corollary}{Corollary}
\newtheorem*{corollary*}{Corollary}
\newtheorem{proposition}{Proposition}
\newtheorem*{proposition*}{Proposition}
\newtheorem{conjecture}{Conjecture}
\newtheorem*{conjecture*}{Conjecture}
{\theoremstyle{definition} \newtheorem{definition}{Definition}
\newtheorem*{definition*}{Definition}
\newtheorem{example}{Example}
\newtheorem*{example*}{Example}
\newtheorem{remark}{Remark}
\newtheorem*{remark*}{Remark}
\newtheorem{note}{Note}
\newtheorem*{note*}{Note}
}

\makeatother

\let\Cal\mathcal


\makeatletter
\newcommand{\g}{\frak g}
\newcommand{\hg}{\hat {\frak g} }
\newcommand{\hn}{\hat {\frak n} }
\newcommand{\h}{\frak h}
\newcommand{\V}{\Cal V}
\newcommand{\hh}{\hat {\frak h} }
\newcommand{\n}{\frak n}
\newcommand{\Z}{\Bbb Z}
\newcommand{\N}{{\Bbb Z} _{> 0} }
\newcommand{\Zp} {\Z _ {\ge 0} }
\newcommand{\C}{\Bbb C}
\newcommand{\Q}{\Bbb Q}
\newcommand{\1}{\bf 1}
\newcommand{\wt}{{\rm {wt} }   }
\makeatother
\begin{document}
\begin{center}\Large For every positive rank ell, the standard quadratic bosonic embedding of ell copies of level-minus-one affine sl2 in two level-minus-one-half affine sp(2ell) vacuum algebras has coset equal to the global-sign orbifold of the rank-ell Heisenberg vertex algebra\end{center}
\noindent\textbf{Stable ID:} 2143\par
\noindent\textbf{Authors:} Drazen Adamovic; Ozren Perse\par
\noindent\textbf{Original work:} The Vertex Algebra M(1)+\allowbreak{} and Certain Affine Vertex Algebras of Level -1\par
\noindent\textbf{Version:} Official SIGMA 8 (2012) article 040, DOI 10.3842/SIGMA.2012.040; journal PDF and native TeX acquired 2026-10-01, exact bytes pinned by SHA256\par
\noindent\textbf{Selected task scope:} Full Theorem2 for ell>=\allowbreak{}1, over C, in the exact Section4-5 Weyl realization: U=\allowbreak{}L\_(A1\textasciicircum{}(1))(-Lambda0) tensor-power ell is generated by the quadratic e(i),f(i),h(i) of(9) inside L\_(C\_ell\textasciicircum{}(1))(-Lambda0/2) tensor L\_(C\_ell\textasciicircum{}(1))(-Lambda0/2). Then Com(U,that tensor VOA)=\allowbreak{}M\_h(1)\textasciicircum{}+\allowbreak{}, where h is spanned by the ell displayed H(i), pairing -2delta\_ij, and plus means fixed points under the simultaneous global sign H(i)->-H(i). The precise embedding is part of the statement; no arbitrary embedding, positive-level rational coset or separate rank-one/type-C variant is asserted.\par
\noindent\textbf{Difficulty:} H2: The argument identifies the commutant for every rank through an explicit involution-adapted Weyl-mode embedding and a full Virasoro decomposition with negative integral Heisenberg charges. Its substantive exclusion of additional coset states combines positive vacuum grading with that charge decomposition, then restricts to the simultaneous global-sign even sector. The complete free-boson construction and arbitrary-rank coset identification exceed an isolated rank-one Wick-contraction exercise.\par
\noindent\textbf{Editorial boundary note (not author text):} The full original all-positive-rank Theorem2 is retained with the exact quadratic embedding, Weyl-mode/involution construction, tensor even-sector embedding, current generators, Heisenberg inner product, full Virasoro splitting and Proposition1 no-extra-state proof. The involution negates all Heisenberg generators simultaneously, rather than taking independent rank-one orbifolds. Standard Feingold-Frenkel and VOA/Sugawara foundations remain cited prerequisites; bounded finite Wick contractions follow the complete displayed quadratic fields. No rank-one, type-C, module-classification or later branching result is separately selected.\par
\noindent\textbf{Source:} \url{https://sigma-journal.com/2012/040/}\quad DOI: \url{https://doi.org/10.3842/SIGMA.2012.040}\par
\noindent\textbf{License and attribution:} Copyright retained by the named authors. Published by SIGMA under CC BY 4.0: \url{https://creativecommons.org/licenses/by/4.0/}. Source: \url{https://sigma-journal.com/about.html}. Adaptation consists of standalone excerpt formatting, cover and numbering restoration; original author language and mathematical text are retained.\par
\noindent\textbf{Typesetting adaptation (not author text):} The legacy Cal command in the author macros is provided as the standard mathcal alias for portable standalone TeX; original mathematical author text and fonts are retained.\par
\medskip\hrule\medskip\noindent\textbf{Author text begins below.}\par
\setcounter{section}{1}
\setcounter{subsection}{0}
\setcounter{subsubsection}{0}
\setcounter{equation}{2}
% Original source lines 205--624
\section{Preliminaries}

Let $V$ be a vertex algebra \cite{B, FHL, FLM,LL}. For a subalgebra $U$ of $V$ denote by
\[
{\rm Com}(U,V)= \{ v \in V \; | \; u_n v =0 \mbox{ for all } u \in U,\;
n \geq 0 \} \]
%
 the commutant of $U$ in $V$ (cf.~\cite{FZ, GKO,LL}). Then, ${\rm Com}(U,V)$ is a
subalgebra of $V$ (also called coset vertex algebra).



Let ${\frak h}$ be a f\/inite-dimensional vector space equipped with a
nondegenerate symmetric bilinear form $\langle \cdot, \cdot
\rangle$, considered as an Abelian Lie algebra. Let $\hat{\frak h} =
{\frak h} \otimes {\C}[t,t^{-1}] \oplus {\C} K$ be its af\/f\/inization
with the center $K$. Then the free bosonic Fock space
$M(1) = S({\frak h} \otimes t^{-1} \C[t^{-1}])$ is a simple vertex
operator algebra of central charge $\ell = \dim {\frak h}$, with
Virasoro vector
\[
 \omega =\frac{1}{2} \sum _{i=1}^{\ell} h^{(i)}(-1)^2 {\1},\]
  where
$\{ h^{(1)}, \ldots , h^{(\ell)} \} $ is any orthonormal basis of
${\frak h}$ (cf.~\cite{FLM, LL}). We shall also use the
notation $M_{\frak h}(1)$ to emphasize the associated vector space
${\frak h}$.

Vertex algebra $M(1)$  has an order 2 automorphism which is lifted
from the map $h \mapsto -h$, for $h \in {\frak h}$. Denote by  $M(1)
^+$ (or $M_{\frak h}(1) ^+$) the subalgebra of invariants of that
automorphism. The irreducible modules for $M(1) ^+$ were classif\/ied
in \cite{DN1} and \cite{DN2}. For $\ell =1$, it was proved in
\cite{DG} that $M(1) ^+$ is generated by $\omega$ and one primary
vector of conformal weight $4$, so it is isomorphic to a
$W(2,4)$-algebra with central charge $1$.

Let $\frak g$ be the simple Lie algebra of type $X_{n}$, and
$\hat{\frak g}$ the associated af\/f\/ine Lie algebra of type
$X_{n}^{(1)}$. For any weight $\Lambda$ of $\hat{\frak g}$, denote
by $L_{X_{n} ^{(1)}}(\Lambda)$ the irreducible highest weight
$\hat{\frak g}$-module. Denote by~$\Lambda _i$, $i=0, \ldots, n$ the
fundamental weights of $\hat{\frak g}$ (cf.~\cite{K1}). We shall
also use the notation $V_{X_n}(\mu)$ for a~highest weight $\frak
g$-module of highest weight~$\mu$, and $\omega _i$, $i=1, \ldots, n$
for the fundamental weights of~$\frak g$.


{\sloppy For any $k \in \C$, denote by $N_{X_{n} ^{(1)}}(k \Lambda _0)$ the
generalized Verma $\hat{\frak g}$-module with highest weight~$k
\Lambda _0$. Then, $N_{X_{n} ^{(1)}}(k \Lambda _0)$ is a vertex
operator algebra of central charge $\frac{k \dim \frak
g}{k+h^{\vee}}$, for any $k \neq - h^{\vee}$, with Virasoro vector
obtained by Sugawara construction:
\begin{gather} \label{Vir-Sugawara}
\omega=\frac{1}{2(k+h^{\vee})}\sum_{i=1}^{\dim {\frak g}}
a^{i}(-1)b^{i}(-1){\bf 1},
\end{gather}
where $\{a^{i}\}_{i=1, \dots, \dim {\frak g}}$ is an arbitrary basis
of ${\frak g}$, and $\{b^{i}\}_{i=1, \dots, \dim {\frak g}}$  the
corresponding dual basis of~${\frak g}$ with respect to the
symmetric invariant bilinear form, normalized by the condition that
the length of the highest root is~$\sqrt{2}$ (cf.~\cite{FrB,FZ,K,LL,L}).

}

It follows that any quotient of $N_{X_{n} ^{(1)}}(k \Lambda _0)$ is
a vertex operator algebra, for $k \neq - h^{\vee}$. Specially,
$L_{X_{n} ^{(1)}}(k \Lambda _0)$ is a simple vertex operator
algebra, for any $k \neq - h^{\vee}$.



\section[Simple Lie algebras of type $C_{\ell}$ and $A_{2\ell -1}$]{Simple Lie algebras of type $\boldsymbol{C_{\ell}}$ and $\boldsymbol{A_{2\ell -1}}$}

Consider two $2 \ell$-dimensional vector spaces $A_1 = \oplus
_{i=1} ^{2 \ell}{\C} a_i ^+$, $ A_2 = \oplus _{i=1}^ {2 \ell} {\C}
a_i^-$ and let $A=A_1 {\oplus} A_2$. The Weyl algebra $W_{2 \ell}$ is
the complex associative algebra generated by $A$ with non-trivial
relations
\[
[a_i ^+,a_j^-]=\delta_{i,j} , \qquad 1 \le i,j \le 2 \ell.
\]
%
The normal ordering on $A$ is def\/ined by
\[
:\!xy\!:    = \frac 1 2 (xy+yx),  \qquad x,y \in A .
\]
Def\/ine
\[
e_{\epsilon_i-\epsilon_j}^A=  :\!a_i ^+ a_j^-\!:, \qquad
f_{\epsilon_i-\epsilon_j}^A=  :\!a_j ^+ a_i^-\!:, \qquad 1 \le i,j
\le 2 \ell, \quad i<j,
\]
and
\[
H_i=-:\!a_i ^+ a_i^-\!:,  \qquad 1 \le i \le 2 \ell.
\]
Then the Lie algebra $\frak g _1$ generated by the set
\[
\big\{ e_{\epsilon_i-\epsilon_j}^A, f_{\epsilon_i-\epsilon_j}^A \; \vert
\; 1 \le i,j \le 2 \ell, \ i < j \big\}
\]
is the simple Lie algebra of type $A_{2 \ell -1}$ (cf.~\cite{Bou}
and~\cite{FF}). The Cartan subalgebra $\frak h _1$ is spanned by
\[
\{ H_i - H_{i+1}  \; \vert \; 1 \le i \le 2 \ell -1 \}.
\]


Let $\theta$ be the automorphism of $W_{2 \ell}$ of order two given
by
%
\[  a_{i} ^+ \mapsto  a_{2\ell+1-i} ^{-}, \qquad a_{2\ell+1-i} ^-
\mapsto  a_{i} ^{+}, \qquad a_{i} ^-  \mapsto - a_{2\ell+1-i} ^+,
\qquad a_{2\ell+1-i} ^+ \mapsto - a_{i} ^-,\]
for $i =1, \dots, \ell$. Clearly, $\frak g _1$ is $\theta$-stable
and \[
\theta (e_{\epsilon_i-\epsilon_j}^A)=- e_{\epsilon_{2\ell
+1-j}-\epsilon_{2\ell +1-i}}^A, \qquad \theta
(e_{\epsilon_i-\epsilon_{2\ell +1-j}}^A)=
e_{\epsilon_{j}-\epsilon_{2\ell +1-i}}^A,
\]
%
and similarly for root vectors associated to negative roots.


The subalgebra $\frak g$ of $\frak g _1$ generated by
\begin{gather*}
  e_{2 \epsilon _{i}} =a_i ^+  a_{2\ell +1-i} ^-, \qquad
    f_{2 \epsilon _{i}}=a_i ^-  a_{2\ell +1-i} ^+,
  \\
  e_{\epsilon _{i}+ \epsilon _{j}}=\frac{1}{2}(a_i ^+
 a_{2\ell +1-j} ^- + a_j ^+  a_{2\ell +1-i} ^-), \qquad
 f_{\epsilon _{i}+ \epsilon _{j}}=\frac{1}{2}(a_i ^- a_{2\ell +1-j}
^+ + a_j ^-  a_{2\ell +1-i} ^+),   \\
  e_{\epsilon _{i}- \epsilon _{j}}=\frac{1}{2}(a_i ^+
 a_{j} ^- - a_{2\ell +1-j} ^+
 a_{2\ell +1-i} ^-), \qquad
f_{\epsilon _{i}- \epsilon _{j}}=\frac{1}{2}(a_j ^+
 a_{i} ^- - a_{2\ell +1-i} ^+
 a_{2\ell +1-j} ^-),
\end{gather*}
for $i,j =1, \dots, \ell$, $i<j$, is the simple Lie algebra of type
$C_{\ell}$. The Cartan subalgebra $\frak h$ is spanned by
\[
\{ H_i - H_{2\ell +1-i}  \; \vert \; 1 \le i \le  \ell \}.
\]
Clearly, $\theta$ acts as $1$ on $\frak g$. Furthermore, \[
e_{\epsilon _{1}+ \epsilon _{2}} ^*=\frac{1}{2}(a_1 ^+
 a_{2\ell -1} ^- - a_2 ^+  a_{2\ell} ^-) \in \frak g _1
\]
is a highest weight vector for $\frak g$, which generates the
irreducible $\frak g$-module $V_{C_{\ell}}(\omega _2)$. Clearly,
$\theta$~acts as $-1$ on $V_{C_{\ell}}(\omega _2)$. We obtain the
decomposition
\begin{gather} \label{decomp-fin1}
 \frak g _1 \cong \frak g \oplus V_{C_{\ell}}(\omega _2).
  \end{gather}
Since
%
\[ \dim V_{A_{2\ell -1}} (n \omega_1) = \dim V_{C_{\ell }} (n
\omega_1),\] one easily concludes that the irreducible $\frak g
_1$-module $V_{A_{2\ell -1}} (n \omega_1)$ remains irreducible when
restricted to $\frak g$. Thus,
\begin{gather} \label{decomp-fin2} V_{A_{2\ell -1} } (n \omega_1) \cong V_{C_{
\ell  } } (n \omega_1) \qquad \mbox{for} \quad n \in \Zp.
\end{gather}
Similarly
\begin{gather} \label{decomp-fin3} V_{A_{2\ell -1} } (n \omega_{2 \ell -1})
\cong V_{C_{ \ell  } } (n \omega_1) \qquad \mbox{for} \quad n \in \Zp.
\end{gather}
We will consider certain af\/f\/ine analogues of relations
(\ref{decomp-fin1})--(\ref{decomp-fin3}).

In what follows we shall need the following decompositions of $\frak
g$-modules:
%
\begin{gather}
V_{C_{\ell }} (\omega_2) \otimes V_{C_{\ell }} ( \omega_2) \cong
V_{C_{\ell }} (2 \omega_2)\oplus V_{C_{\ell }} ( \omega_1 +
\omega_3) \oplus V_{C_{\ell }} (\omega_4)
\nonumber \\
\hphantom{V_{C_{\ell }} (\omega_2) \otimes V_{C_{\ell }} ( \omega_2) \cong}{}
 \oplus V_{C_{\ell }} (2\omega_1) \oplus
V_{C_{\ell}} ( \omega_2) \oplus V_{C_{\ell }} (0)  \quad ( \ell \ge 4),  \nonumber \\
 V_{C_{3}} (\omega_2) \otimes V_{C_{3}} ( \omega_2) \cong
 V_{C_{3}} (2 \omega_2)\oplus V_{C_{3}} ( \omega_1 +
\omega_3) \oplus V_{C_{3}} (2\omega_1) \oplus
V_{C_{3}} ( \omega_2) \oplus V_{C_{3}} (0), \nonumber \\
 V_{C_{2}}(\omega _2) \otimes V_{C_{2}}(\omega _2) \cong
V_{C_{2}}(2 \omega _2) \oplus V_{C_{2}}(2 \omega _1) \oplus V_{C_{2}}(0). \label{tens-pr-decomp}
 \end{gather}



\section{Weyl vertex algebras and symplectic af\/f\/ine Lie algebras}

The Weyl algebra $W_{\ell}( \tfrac{1}{2} + {\Z})$ is a complex
associative algebra generated by
\[  a ^{\pm} _{i}(r),  \qquad   r  \in \tfrac{1}{2} + {\Z}, \quad 1 \le i
\le \ell
\] with non-trivial relations
\begin{gather*}
 [a_{i} ^{+}(r)  , a_{j} ^{-}(s)  ]  =\delta_{r+s,0} \delta_{i,j},
\end{gather*}
where $r, s \in {\tfrac{1}{2}}+ {\Z}$, $ i, j \in \{1, \dots, \ell
\}$.



 Let $M_{\ell}$ be the irreducible $W_{\ell}( \tfrac{1}{2} + {\Z})$-module generated by
 the
cyclic vector ${\1}$ such that
\[ a^{\pm} _i  (r) {\1}     =  0 \qquad \mbox{for} \quad r > 0 ,  \quad 1
\le i \le \ell.
\]



Def\/ine the following   f\/ields on $M_{\ell}$
%
\[   a_i ^{\pm }(z) = \sum_{ n \in   {\Z}
 } a_i ^{\pm }(n+ \tfrac{1}{2} )  z ^{-n- 1}.
 \]
 The f\/ields $a_{i} ^{\pm }(z)$, $i =1, \dots, \ell$ generate on $M_{\ell}$  the
unique structure of a simple vertex algebra (cf.~\cite{FrB, K}). Let us denote the corresponding vertex operator by $Y$.

We have the following Virasoro vector in $M_{\ell}$:
\begin{gather} \label{vir-bos}  \omega=\frac{1}{2} \sum_{i=1} ^{\ell} \big( a_i
^-  (- \tfrac{3}{2} ) a_i ^{+}  (-\tfrac{1}{2}  ) -a_i ^+  (-  \tfrac{3}{2} )
a_i ^{-} (-\tfrac{1}{2}  ) \big) {\1}.
\end{gather}
Let $Y(\omega,z) = \sum\limits_{n \in {\Z} } L(n) z ^{-n-2}$.
Then $M_{\ell}$ is $\tfrac{1}{2}{\Zp}$-graded with respect to
$L(0)$:
\[ M_{\ell} := \bigoplus_{ m \in \tfrac{1}{2} \Zp} M_{\ell} (m),
\qquad M_{\ell} (m)= \{ v \in M_{\ell} \; \vert \;  L(0) v = m v \}.\]
Note that $M_{\ell} (0) = {\C} {\1}$. For $ v \in M_{\ell} (m)$ we
shall write $\wt (v) = m$.

The following result is well-known.

\begin{theorem}[\cite{FF}] \label{ff}
We have
\[  M_{\ell} \cong L_{C_{\ell} ^{(1)}}( -\tfrac{1}{2}\Lambda _{0})
\bigoplus L_{C_{\ell} ^{(1)}}( -\tfrac{3}{2}\Lambda _{0} +
\Lambda_1). \]
In the case $\ell =1$ we have \[  M_1 \cong L_{A_1 ^{(1)}}(
-\tfrac{1}{2}\Lambda _{0}) \bigoplus L_{A_1 ^{(1)}}(
-\tfrac{3}{2}\Lambda _{0} + \Lambda_1). \]
\end{theorem}

\begin{remark} The highest weights of modules from Theorem \ref{ff}
are admissible in the sense of~\cite{KW}. Representations of vertex
operator algebras associated to af\/f\/ine Lie algebras of type~$A_1
^{(1)}$ and~$C_{\ell} ^{(1)}$ with admissible highest weights were
studied in~\cite{AM} and~\cite{A-1994}.
\end{remark}

We shall now consider the vertex algebra $M_{2 \ell}$ and its
subalgebra $L_{C_{\ell} ^{(1)}}( -\tfrac{1}{2}\Lambda _{0}) \otimes
L_{C_{\ell} ^{(1)}}( -\tfrac{1}{2}\Lambda _{0})$. For $\ell =1$,
$L_{A_1 ^{(1)}}(-\tfrac{1}{2}\Lambda_0) \otimes L_{A_1
^{(1)}}(-\tfrac{1}{2}\Lambda_0)$ is a subalgebra of $M_{2}$.





Let $\theta : M_{2 \ell} \rightarrow M_{2 \ell}$ be the automorphism
of order two of  the vertex algebra $M_{2 \ell}$ which is lifted
from the following automorphism of the Weyl algebra
\begin{alignat*}{3}
& a_{i} ^+ (s) \mapsto  a_{2\ell+1-i} ^{-} (s), \qquad
&& a_{2\ell+1-i} ^- (s) \mapsto  a_{i} ^{+} (s), &  \\
& a_{i} ^-(s) \mapsto - a_{2\ell+1-i} ^+ (s), \qquad && a_{2\ell+1-i}
^+(s) \mapsto - a_{i} ^- (s), &
 \end{alignat*}
for $i =1, \dots, \ell$ and  $s \in \tfrac{1}{2} + {\Z}$.

If we have a subalgebra $U \subset M_{2 \ell}$ which is
$\theta$-stable, we def\/ine
\[
U ^0 = \{ u \in U \; \vert \; \theta (u) = u \}, \qquad U ^1 = \{ u
\in U \; \vert \; \theta (u) = -u \}.
\]

Def\/ine the following vectors in $M_{2 \ell}$:
\begin{alignat*}{3}
&b_{i} ^+ = \frac{1}{\sqrt{2}}( a_{i} ^+ (-\tfrac{1}{2}) +
a_{2\ell+1-i} ^- (-\tfrac{1}{2}) ) {\1},\qquad
&& b_{2\ell+1-i} ^+ = \frac{\sqrt{-1}}{\sqrt{2}}( a_{i} ^+ (-\tfrac{1}{2}) - a_{2\ell+1-i} ^- (-\tfrac{1}{2}) ) {\1}, &  \\
& b_{i} ^- = \frac{1}{\sqrt{2}}( a_{i} ^- (-\tfrac{1}{2}) -
a_{2\ell+1-i}^+ (-\tfrac{1}{2}) ) {\1}, \qquad &&b_{2\ell+1-i} ^- =
\frac{-\sqrt{-1}}{\sqrt{2}}( a_{i} ^- (-\tfrac{1}{2}) +
a_{2\ell+1-i} ^+ (-\tfrac{1}{2}) ) {\1}, &
\end{alignat*} for $i =1, \dots, \ell$. Then the subalgebra generated by $b_i
^+$, $b _i ^-$ (resp.\ $b_{2\ell+1-i} ^+$, $b _{2\ell+1-i} ^-$), for $i
=1, \dots, \ell$, is isomorphic to $M_{\ell}$.
Since
\[ \theta (b_i ^{\pm}) = b_i ^{\pm}, \qquad  \theta (b_{2\ell+1-i}
^{\pm}) = - b_{2\ell+1-i} ^{\pm},\]
for $i =1, \dots, \ell$, we have \[ M_{2 \ell}^{0} \cong M_{\ell}
\otimes M_{\ell} ^0 \cong M_{\ell} \otimes L_{C_{\ell} ^{(1)}}(
-\tfrac{1}{2}\Lambda _{0}),\] and for $\ell =1$: \[ M_{2}^{0} \cong
M_{1} \otimes M_{1} ^0 \cong M_{1} \otimes L_{A_1 ^{(1)}}(
-\tfrac{1}{2}\Lambda _{0}).\]




\section[Commutant of $L_{A_1 ^{(1)}}(- \Lambda_0) ^{\otimes \ell}$ in  $L_{C_{\ell} ^{(1)}}(-\frac{1}{2}\Lambda_0)
\otimes L_{C_{\ell} ^{(1)}}(-\frac{1}{2}\Lambda_0)$]{Commutant of $\boldsymbol{L_{A_1 ^{(1)}}(- \Lambda_0) ^{\otimes \ell}}$ in  $\boldsymbol{L_{C_{\ell} ^{(1)}}(-\frac{1}{2}\Lambda_0)
\otimes L_{C_{\ell} ^{(1)}}(-\frac{1}{2}\Lambda_0)}$}

In this section we use Weyl vertex algebra $M_{2 \ell}$ to study
certain subalgebras of $L_{C_{\ell} ^{(1)}}(-\tfrac{1}{2}\Lambda_0)
\otimes L_{C_{\ell} ^{(1)}}(-\tfrac{1}{2}\Lambda_0)$. We determine
the commutants
\[
{\rm Com} \big( L_{A_1 ^{(1)}}(- \Lambda_0) ^{\otimes
\ell},L_{C_{\ell} ^{(1)}}(-\tfrac{1}{2}\Lambda_0) \otimes
L_{C_{\ell} ^{(1)}}(-\tfrac{1}{2}\Lambda_0)\big)
\]
and \[
{\rm Com} \big( L_{A_1 ^{(1)}}(-\Lambda_0),L_{A_1
^{(1)}}(-\tfrac{1}{2}\Lambda_0) \otimes L_{A_1
^{(1)}}(-\tfrac{1}{2}\Lambda_0) \big).
\]



Let $U$ be the vertex subalgebra  of $M_{\ell} ^0 \otimes M_{\ell}
^0 \subset M_{2 \ell} ^0$ generated by
\begin{gather}
e^{(i)}=a_i ^+ (-\tfrac{1}{2}) a_{
2\ell +1-i} ^-( -\tfrac{1}{2}) {\1}, \qquad  f^{(i)}=a_ i ^-
(-\tfrac{1}{2}) a_{2\ell +1-i} ^+(
-\tfrac{1}{2}) {\1} \qquad \mbox{and} \nonumber \\
 h^{(i)}= (-a_ i ^+ (-\tfrac{1}{2}) a_i^-( -\tfrac{1}{2}) +
a_{2\ell +1-i} ^+ (-\tfrac{1}{2}) a_{2\ell +1-i} ^-( -\tfrac{1}{2})){\1}, \label{gen-1-higher}
 \end{gather}
$i =1, \dots, \ell$. It is clear that $U$ is isomorphic to the tensor product
\[
 \underbrace{ L_{A_1 ^{(1)}}( -\Lambda _{0}) \otimes \cdots \otimes  L_{A_1 ^{(1)}}( -\Lambda _{0})}_{\ell \ \  \mbox{times} }
\]
of $\ell$ copies of the af\/f\/ine vertex algebra $ L_{A_1 ^{(1)}}(-\Lambda _{0})$.

Def\/ine also
 \[
 H ^{(i)} =   (a_i ^+ (-\tfrac{1}{2}) a_i ^- (-\tfrac{1}{2})+a_{2\ell +1-i} ^+ (-\tfrac{1}{2})
 a_{2\ell +1-i} ^- (-\tfrac{1}{2}) ) {\1}, \qquad i=1, \dots,
 \ell.
 \]
 Then
 \[ H ^{(i)} \in \mbox{Com} ( U, M_{2\ell}), \qquad i=1, \dots,  \ell.
 \]

 Let \[ {\frak h}= \bigoplus_{i =1 } ^{ \ell} {\C} H ^{(i)}.\] Then ${\frak h}$ can be
 considered as an Abelian Lie algebra, and the components of the vertex operators
 \[
 Y(h,z) = \sum_{\n \in {\Z} } h(n) z ^{-n-1}, \qquad h \in {\frak  h},
 \]
 def\/ine a representation of the associated Heisenberg algebra $\widehat{\frak h}$.
 Moreover, ${\frak h}$ generates the subalgebra of the~$M_{2 \ell}$ which is isomorphic to the
 Heisenberg vertex algebra $M_{\frak h} (1)$ with central charge~$\ell$.
 We also note that $\langle H ^{(i)}, H ^{(j)} \rangle = H ^{(i)}(1)H ^{(j)}= -2 \delta _{ij}$.


Using relations (\ref{gen-1-higher}) one can show that the Virasoro
vector (\ref{vir-bos}) (in $M_{2 \ell}$) can be written in a~form:
\[ \omega =  \omega_1   + \omega_2 \]
where \[\omega_1 =\frac{1}{2} \sum_{i = 1 } ^{\ell} \big(e
^{(i)}(-1)f^{(i)}(-1)+f ^{(i)} (-1)e^{(i)}(-1)+ \tfrac{1}{2} h ^{(i)}
(-1)^2\big){\1}
\] is the  Virasoro vector  in $U$  and
\[\omega_2 = -\frac{1}{4}\sum_{i = 1 } ^{\ell} H ^{(i) }(-1) ^{2}
{\1}\] is the Virasoro vector   in $M_{\frak h} (1)$.
For $i=1,2$ let
$Y(\omega_i , z) = \sum\limits_{n\in {\Z} } L_i (n) z ^{-n-2}$.

\begin{proposition} \label{com-bos-higher}
We have:
\[
{\rm Com} (U, M_{2\ell}) = M_{\frak h} (1).
\]
\end{proposition}

\begin{proof}
It is clear that
 $W= \mbox{Com} ( U, M_{2\ell})$ contains a subalgebra isomorphic to $M_{\frak h}
(1)$.

Assume now that $M_{\frak h} (1) \ne W$. Then there is a vector
$ w \in W$, ${\wt} (w) > 0$ such that
\[ H ^{(i)}(n) w =  \delta_{n,0} \lambda _{(i)} w, \qquad n \in {\Zp},
\quad \lambda _{(i)} \in {\Z}, \quad i = 1, \dots, \ell.
\]
(Note that each  $H^{(i)}(0)$ acts semisimply on $M_{2\ell}$ with
eigenvalues in ${\Z}$.)
By the def\/inition of $W$ we have that $L_1 (0) w = 0$. Therefore
\[ L(0) w = L_2(0) w = - \frac{1}{4} \sum_{i = 1 } ^{\ell} (\lambda
_{(i)})  ^2 w. \] This contradicts the fact that $\wt (w) > 0$. So,
$W= M_{\frak h} (1)$.
\end{proof}


Since $\theta (H^{(i)}) = -H^{(i)}$ for $i=1, \dots, \ell$ we have
that $M_{\frak h} (1)^+ =M_{\frak h}(1) ^0$ is a subalgebra of
$M_{\ell} ^0 \otimes M_{\ell} ^0$.
 Therefore the  vertex algebra $L_{C_{\ell} ^{(1)}}(-\tfrac{1}{2}\Lambda_0) \otimes L_{C_{\ell} ^{(1)}}(-\tfrac{1}{2}\Lambda_0)$ contains a subalgebra isomorphic
to $U \otimes M_{\frak h} (1) ^+$. By using Proposition~\ref{com-bos-higher} we obtain the following theorem.

\begin{theorem} \label{coset-bos-r-higher}We have:
\[ {\rm Com} \big( L_{A_1 ^{(1)}}(- \Lambda_0) ^{\otimes
\ell},L_{C_{\ell} ^{(1)}}(-\tfrac{1}{2}\Lambda_0) \otimes
L_{C_{\ell} ^{(1)}}(-\tfrac{1}{2}\Lambda_0) \big) \cong M_{\frak h} (1)
^+.\]
\end{theorem}
\begin{thebibliography}{99}
\bibitem[2]{A-1994}
Adamovi{\'c} D., Some rational vertex algebras, \textit{Glas. Mat. Ser.~III}
  \textbf{29(49)} (1994), 25--40, \href{http://arxiv.org/abs/q-alg/9502015}{q-alg/9502015}.

\bibitem[5]{AM}
Adamovi{\'c} D., Milas A., Vertex operator algebras associated to modular
  invariant representations for {$A^{(1)}_1$}, \textit{Math. Res. Lett.}
  \textbf{2} (1995), 563--575, \href{http://arxiv.org/abs/q-alg/9509025}{q-alg/9509025}.

\bibitem[10]{B}
Borcherds R.E., Vertex algebras, {K}ac--{M}oody algebras, and the {M}onster,
  \href{http://dx.doi.org/10.1073/pnas.83.10.3068}{\textit{Proc. Nat. Acad. Sci. USA}} \textbf{83} (1986), 3068--3071.

\bibitem[11]{Bou}
Bourbaki N., \'El\'ements de math\'ematique. Fasc.~XXXVIII: Groupes et
  alg\`ebres de Lie, \textit{Actualit\'es Scientifiques et Industrielles}, Vol.
  1364, Hermann, Paris, 1975.

\bibitem[12]{DG}
Dong C., Griess R.L., Rank one lattice type vertex operator algebras and their
  automorphism groups, \href{http://dx.doi.org/10.1006/jabr.1998.7498}{\textit{J.~Algebra}} \textbf{208} (1998), 262--275,
  \href{http://arxiv.org/abs/q-alg/9710017}{q-alg/9710017}.

\bibitem[15]{DN1}
Dong C., Nagatomo K., Classif\/ication of irreducible modules for the vertex
  operator algebra {$M(1)^+$}, \href{http://dx.doi.org/10.1006/jabr.1998.7784}{\textit{J.~Algebra}} \textbf{216} (1999),
  384--404, \href{http://arxiv.org/abs/math.QA/9806051}{math.QA/9806051}.

\bibitem[16]{DN2}
Dong C., Nagatomo K., Classif\/ication of irreducible modules for the vertex
  operator algebra {$M(1)^+$}. {II}.~{H}igher rank, \href{http://dx.doi.org/10.1006/jabr.2000.8716}{\textit{J.~Algebra}}
  \textbf{240} (2001), 289--325, \href{http://arxiv.org/abs/math.QA/9905064}{math.QA/9905064}.

\bibitem[17]{FF}
Feingold A.J., Frenkel I.B., Classical af\/f\/ine algebras, \href{http://dx.doi.org/10.1016/0001-8708(85)90027-1}{\textit{Adv. Math.}}
  \textbf{56} (1985), 117--172.

\bibitem[18]{FrB}
Frenkel E., Ben-Zvi D., Vertex algebras and algebraic curves,
  \textit{Mathematical Surveys and Monographs}, Vol.~88, American Mathematical
  Society, Providence, RI, 2001.

\bibitem[19]{FHL}
Frenkel I.B., Huang Y.Z., Lepowsky J., On axiomatic approaches to vertex
  operator algebras and modules, \textit{Mem. Amer. Math. Soc.} \textbf{104}
  (1993), no.~494.

\bibitem[20]{FLM}
Frenkel I.B., Lepowsky J., Meurman A., Vertex operator algebras and the
  {M}onster, \textit{Pure and Applied Mathematics}, Vol.~134, Academic Press
  Inc., Boston, MA, 1988.


\bibitem[21]{FZ}
Frenkel I.B., Zhu Y., Vertex operator algebras associated to representations of
  af\/f\/ine and {V}irasoro algebras, \href{http://dx.doi.org/10.1215/S0012-7094-92-06604-X}{\textit{Duke Math.~J.}} \textbf{66} (1992),
  123--168.

\bibitem[22]{GKO}
Goddard P., Kent A., Olive D., Virasoro algebras and coset space models,
  \href{http://dx.doi.org/10.1016/0370-2693(85)91145-1}{\textit{Phys. Lett.~B}} \textbf{152} (1985), 88--92.

\bibitem[23]{K1}
Kac V.G., Inf\/inite-dimensional {L}ie algebras, 3rd ed., \href{http://dx.doi.org/10.1017/CBO9780511626234}{Cambridge University
  Press}, Cambridge, 1990.

\bibitem[24]{K}
Kac V.G., Vertex algebras for beginners, \textit{University Lecture Series},
  Vol.~10, 2nd ed., American Mathematical Society, Providence, RI, 1998.


\bibitem[25]{KW}
Kac V.G., Wakimoto M., Modular invariant representations of
  inf\/inite-dimensional {L}ie algebras and superalgebras, \href{http://dx.doi.org/10.1073/pnas.85.14.4956}{\textit{Proc. Nat.
  Acad. Sci. USA}} \textbf{85} (1988), 4956--4960.

\bibitem[27]{LL}
Lepowsky J., Li H., Introduction to vertex operator algebras and their
  representations, \href{http://dx.doi.org/10.1007/978-0-8176-8186-9}{\textit{Progress in Mathematics}}, Vol.~227, Birkh\"auser
  Boston Inc., Boston, MA, 2004.

\bibitem[28]{L}
Li H.S., Local systems of vertex operators, vertex superalgebras and modules,
  \href{http://dx.doi.org/10.1016/0022-4049(95)00079-8}{\textit{J.~Pure Appl. Algebra}} \textbf{109} (1996), 143--195,
  \href{http://arxiv.org/abs/hep-th/9406185}{hep-th/9406185}.

\end{thebibliography}
\end{document}
